%%%%%%%%%%%%%%%%%%%%%%%%%%%%%%%%%%%%%%%%%%%%%%%%%%%%%%%%%%%%%%%%%%%%%%
% Lớp 10 | Véctơ (chưa xét tọa độ) | Tích vô hướng (chưa xét tọa độ)
% Dạng 0H5?4-5: Các bài toán tìm điểm và tập hợp điểm
% Số câu: 1  (sinh tự động bởi tach_cau.py; từ nay thầy sửa trực tiếp file này)
%%%%%%%%%%%%%%%%%%%%%%%%%%%%%%%%%%%%%%%%%%%%%%%%%%%%%%%%%%%%%%%%%%%%%%

%%%%%%%%%%%%%%%% TRẢ LỜI NGẮN %%%%%%%%%%%%%%%%

% nguồn: MaHoaNH_TraLoiNgan_Lop10
\begin{bt}%[0H5H4-5]%[ma:021b4e00]
	Cho $\triangle ABC$ đều có cạnh bằng $3$. Điểm $M$ thỏa mãn $MA^2+MB^2+MC^2=15$, khi đó tập hợp  điểm $M$ thuộc đường tròn có bán kính bằng bao nhiêu (làm tròn kết quả đến hàng phần trăm)?
	 \shortans{1,41}
	\loigiai{ 
		Gọi $G$ là trọng tâm tam  giác $ABC$, ta có $GA=GB=GC=\sqrt{3}$ và $\overrightarrow{GA}+\overrightarrow{GB}+\overrightarrow{GC}=\overrightarrow{0}$.\\
		Khi đó, \begin{align*}
		 MA^2+MB^2+MC^2&=\left(\overrightarrow{MG}+\overrightarrow{GA}\right)^2+\left(\overrightarrow{MG}+\overrightarrow{GB}\right)^2+\left(\overrightarrow{MG}+\overrightarrow{GC}\right)^2\\
		 &=3MG^2+GA^2+GB^2+GC^2\\
		&=3MG^2+9. 
	\end{align*}
	Suy ra $15=3MG^2+9\Leftrightarrow MG^2=2\Leftrightarrow MG=\sqrt{2}$.\\ 
	Vậy tập hợp các điểm $M$ là đường tròn tâm $G$ bán  kính $R=\sqrt{2}\approx 1{,}41$.
} 
\end{bt}
